\chapter{Advanced Presentations}

This chapter expands on the \texttt{presentation} doctype overview from
Chapter~\ref{ch:visual} with advanced customization patterns, internal
details, and practical techniques.

\section{Design Philosophy}

OmniLaTeX presentations use KOMA-Script (\texttt{scrartcl}) rather than the
\texttt{beamer} class.  This gives several advantages:

\begin{itemize}
    \item Consistent font, colour, and bibliography handling with all other
          doctypes.
    \item No separate theme system to learn --- institution configs apply
          directly.
    \item Full access to the entire OmniLaTeX macro library, including
          \texttt{minted} code listings, glossaries, and \texttt{tcolorbox}
          environments.
    \item 16:10 aspect ratio (254\,mm $\times$ 190.5\,mm) via direct
          \texttt{\textbackslash paperwidth} assignment, avoiding geometry
          conflicts.
\end{itemize}

The core engine lives in
\texttt{lib/layout/omnilatex-presentation.sty} (v1.13.0).

\section{Colour Customization}

The header, footer, and progress bar use two colours:
\texttt{universityprimary} and \texttt{universitysecondary}.  When no
institution config is loaded, the module defines fallback values:

\begin{minted}{latex}
% Fallback defaults (RGB 0,63,114 and RGB 0,113,188)
% Override in document-settings.sty:
\definecolor{universityprimary}{RGB}{102,0,153}
\definecolor{universitysecondary}{RGB}{153,51,204}
\end{minted}

Institution configs set these colours automatically.  For example, TUHH
defines a dark navy primary and a teal secondary.  The presentation module
checks with \texttt{\textbackslash @ifundefined\{color@universityprimary\}}
before applying its fallback, so any definition in the preamble or an
institution config takes precedence.

\section{Title Page Internals}

The \texttt{\textbackslash maketitle} override fills the upper 35\% of the
page with \texttt{universityprimary}, adds a 2\% accent stripe of
\texttt{universitysecondary}, and centres the title in 28\,pt bold white.
The author appears below in large white, and the institute and date are
placed beneath the coloured region.

\begin{minted}{latex}
\title{My Research Talk}
\author{Jane Doe}
% Institute is not a standard KOMA command; define it manually:
\makeatletter
\gdef\@institute{Department of Computer Science, University of Examples}
\makeatother
\date{June 2026}

\begin{document}
\maketitle
% The title page is followed by a \clearpage
\end{document}
\end{minted}

The page style is set to \texttt{empty} for the title page, so no header or
footer decorations appear.

\section{Frame Environment Details}

Each \texttt{presentationframe} performs these steps in order:

\begin{enumerate}
    \item \texttt{\textbackslash clearpage} --- flush previous content.
    \item \texttt{\textbackslash presentationframemark} --- increment
          \texttt{framecounter}.
    \item \texttt{\textbackslash presentationHeader} --- draw the TikZ header
          bar with title and author.
    \item User content is rendered.
    \item \texttt{\textbackslash presentationFooter} --- draw the footer rule,
          short title, and frame counter.
    \item \texttt{\textbackslash presentationProgressBar} --- draw the 2\,pt
          progress bar.
    \item \texttt{\textbackslash clearpage} --- finalize the slide.
\end{enumerate}

The environment accepts an optional argument (currently unused but reserved
for future per-frame configuration):

\begin{minted}{latex}
\begin{presentationframe}
    \slidetitle{Slide Title}
    Content goes here.
\end{presentationframe}
\end{minted}

\section{Progress Bar Mechanics}

The progress bar requires two-pass compilation to know the total frame count.
Place \texttt{\textbackslash presentationcountframes} at the end of the
document body:

\begin{minted}{latex}
\begin{document}
\maketitle
\begin{presentationframe}
    \slidetitle{First}
    Content.
\end{presentationframe}
\begin{presentationframe}
    \slidetitle{Second}
    Content.
\end{presentationframe}
\presentationcountframes  % <-- snapshots total, resets counter
\end{document}
\end{minted}

Internally, this command copies \texttt{framecounter} into
\texttt{totalframes} and resets \texttt{framecounter} to zero.  On the next
compilation pass, each frame increments \texttt{framecounter} again, and the
progress bar fill width is computed as:

\begin{minted}{latex}
\paperwidth * \ratio{framecounter}{totalframes}
\end{minted}

If \texttt{\textbackslash presentationcountframes} is omitted, the progress
bar shows 0\% fill on every slide (since \texttt{totalframes} remains zero).

\section{Disabling UI Elements}

Two toggle pairs control optional slide decorations:

\begin{minted}{latex}
% Hide navigation symbols (default: off)
\presentationnavsymbolsoff

% Show navigation symbols
\presentationnavsymbolson

% Hide progress bar (default: on)
\presentationprogressoff

% Show progress bar
\presentationprogresson
\end{minted}

Place these commands in the preamble after
\texttt{\textbackslash begin\{document\}} or in
\texttt{config/document-settings.sty}.

\section{Block Customization via tcolorbox}

All five block environments are standard \texttt{tcolorbox} definitions with
the \texttt{enhanced} skin.  The optional first argument passes through to
tcolorbox, allowing per-block overrides:

\begin{minted}{latex}
% Custom width and shadow for one block
\begin{presentationblock}[width=0.6\textwidth, drop shadow]{Tip}
    This block is narrower and has a drop shadow.
\end{presentationblock}

% Override the border colour of a definition block
\begin{definitionblock}[borderline west={3pt}{0pt}{purple}]{Lemma}
    A partial order is reflexive, antisymmetric, and transitive.
\end{definitionblock}
\end{minted}

The five environments and their default styles:

\begin{table}[htbp]
    \centering
    \caption{Block environment defaults}
    \label{tab:block-defaults}
    \begin{tabular}{@{} l l l l @{}}
        \toprule
        \textbf{Environment} & \textbf{Background} & \textbf{Border} & \textbf{Padding} \\
        \midrule
        \texttt{presentationblock} & White        & Primary (2.5\,pt)   & 6\,pt/4\,pt \\
        \texttt{alertblock}        & Red!5        & Red!70!black        & 6\,pt/4\,pt \\
        \texttt{exampleblock}      & Green!5      & Green!50!black      & 6\,pt/4\,pt \\
        \texttt{noteblock}         & Blue!5       & Blue!50!black       & 6\,pt/4\,pt \\
        \texttt{definitionblock}   & Orange!5     & Orange!70!black     & 6\,pt/4\,pt \\
        \bottomrule
    \end{tabular}
\end{table}

\section{Multi-Column Layouts}

The \texttt{presentationcolumns} environment wraps a \texttt{columns[T]}
environment (top-aligned).  Use \texttt{\textbackslash presentationcolumn\{w\}}
to open each column, where \texttt{w} is a fraction of \texttt{textwidth}:

\begin{minted}{latex}
\begin{presentationframe}
    \slidetitle{Two-Column Comparison}
    \begin{presentationcolumns}
        \presentationcolumn{0.45}
            \textbf{Advantages:}
            \begin{itemize}
                \item Reproducible
                \item Version controlled
            \end{itemize}
        \end{column}
        \presentationcolumn{0.45}
            \textbf{Disadvantages:}
            \begin{itemize}
                \item Learning curve
                \item Compile time
            \end{itemize}
        \end{column}
    \end{presentationcolumns}
\end{presentationframe}
\end{minted}

The default column width is 0.48\texttt{textwidth}.  Each column must be
closed explicitly with \texttt{\textbackslash end\{column\}}.

For three-column layouts, adjust widths accordingly:

\begin{minted}{latex}
\begin{presentationcolumns}
    \presentationcolumn{0.30}
        Column 1.
    \end{column}
    \presentationcolumn{0.30}
        Column 2.
    \end{column}
    \presentationcolumn{0.30}
        Column 3.
    \end{column}
\end{presentationcolumns}
\end{minted}

\section{Section Dividers}

The \texttt{\textbackslash presentationSection\{title\}} command creates a
full-width section divider page.  It fills the upper 50\% with
\texttt{universityprimary}, centres the title in 32\,pt bold white, and
includes the footer and progress bar.  The frame counter is incremented
automatically.

\begin{minted}{latex}
\presentationSection{Experimental Setup}
\begin{presentationframe}
    \slidetitle{Equipment}
    We used a calibrated spectrometer\ldots
\end{presentationframe}
\end{minted}

Section dividers count toward the total frame count, so the progress bar
accounts for them correctly.

\section{TikZ Guard Behaviour}

The presentation module auto-loads TikZ if it is not already present:

\begin{itemize}
    \item If \texttt{enabletikz=true} (default), TikZ is loaded before the
          presentation module, so all overlays render.
    \item If \texttt{enabletikz=false}, the module attempts to load TikZ
          itself.  If \texttt{tikz.sty} is unavailable, the
          \texttt{\textbackslash ifpresentation@tikz} flag is set to
          \texttt{false} and all TikZ overlays (header, footer, progress bar,
          title page) are silently skipped.
\end{itemize}

This means presentations always compile even without TikZ, but without any
decorative elements.

\section{Building Presentations}

Use the standard build commands:

\begin{minted}{bash}
# Via Docker
docker run --rm --entrypoint "" \
    -v "$PWD:/workspace" \
    ghcr.io/wyattau/omnilatex:latest \
    python3 build.py --mode prod --force build-example presentation

# Via Nix
nix develop .# --command python3 build.py --mode prod build-example presentation
\end{minted}

The output PDF uses the 16:10 aspect ratio.  For 4:3 projectors, override the
paper dimensions in \texttt{config/document-settings.sty}:

\begin{minted}{latex}
\AtBeginDocument{%
    \setlength{\paperwidth}{254mm}%
    \setlength{\paperheight}{190.5mm}%   % 16:10 (default)
    % \setlength{\paperheight}{190.5mm}%  % uncomment for 4:3: use 254mm x 190.5mm
}
\end{minted}

\section{Common Patterns}

\subsection{Thank-You Slide}

A simple closing slide:

\begin{minted}{latex}
\begin{presentationframe}
    \vfill
    \begin{center}
        {\Huge\bfseries\color{universityprimary}Thank You!}\par
        \vspace{8mm}
        {\Large Questions?}\par
        \vspace{6mm}
        {\normalsize\texttt{jane@example.com}}
    \end{center}
    \vfill
\end{presentationframe}
\end{minted}

\subsection{Listing Slides}

Combine \texttt{minted} with \texttt{presentationblock}:

\begin{minted}{latex}
\begin{presentationframe}
    \slidetitle{Example Implementation}
    \begin{presentationblock}[Python]{Fibonacci}
        \begin{minted}[fontsize=\small]{python}
def fib(n: int) -> int:
    a, b = 0, 1
    for _ in range(n):
        a, b = b, a + b
    return a
        \end{minted}
    \end{presentationblock}
\end{presentationframe}
\end{minted}

\subsection{Figure Slides}

Place figures inside frames without float wrappers:

\begin{minted}{latex}
\begin{presentationframe}
    \slidetitle{System Architecture}
    \begin{center}
        \includegraphics[width=0.75\textwidth]{figures/architecture}
    \end{center}
\end{presentationframe}
\end{minted}

Do not use \texttt{\textbackslash begin\{figure\}} inside frames --- floats
are suppressed by the empty page style.
